\documentclass[11pt]{article}
% Shared preamble for the bite-sized LaTeX doc series (docs/latex/*.tex).
% \input this at the top of each numbered file rather than duplicating the
% package list -- keeps every file's own content close to (and comfortably
% under) 10 pages.
\usepackage[a4paper,margin=2.5cm]{geometry}
\usepackage{amsmath,amssymb}
\usepackage{hyperref}
\usepackage[backend=biber,style=numeric,sorting=none]{biblatex}
\usepackage{listings}
\usepackage{xcolor}
\usepackage{enumitem}
\addbibresource{references.bib}

\lstdefinelanguage{Rust}{
  keywords={fn,let,mut,pub,struct,enum,impl,match,if,else,for,while,return,use,mod,self,Self,const,static},
  keywordstyle=\color{blue}\bfseries,
  comment=[l]{//},
  commentstyle=\color{gray}\itshape,
  string=[b]",
  stringstyle=\color{teal},
  basicstyle=\ttfamily\small,
  showstringspaces=false,
  breaklines=true,
}
\lstset{language=Rust}

\hypersetup{colorlinks=true, linkcolor=blue, citecolor=blue, urlcolor=blue}


\title{The Non-Analytic Critical-Region Term \\ \large Part 2: Numerical Methods}
\author{outram-park-fork-coolprop}
\date{2026-07-10}

\begin{document}
\maketitle

\begin{abstract}
Part 1 introduced the non-analytic term's three auxiliary functions
$\theta$, $\Delta$, $\psi$ and their physical role. This document works
through the first- and second-derivative chain needed for pressure, $c_v$,
$c_p$, and the speed of sound, and the numerical guard required to evaluate
those derivatives \emph{at} the critical point itself, where several
intermediate quantities are formally singular.
\end{abstract}

\section{What derivatives are actually needed}

Every property this crate computes from the Helmholtz surface — pressure,
internal energy, entropy, both heat capacities, the speed of sound — is
built from the five quantities
\[
  \alpha,\quad
  \alpha_\delta = \frac{\partial \alpha}{\partial \delta},\quad
  \alpha_\tau = \frac{\partial \alpha}{\partial \tau},\quad
  \alpha_{\delta\delta} = \frac{\partial^2 \alpha}{\partial \delta^2},\quad
  \alpha_{\tau\tau} = \frac{\partial^2 \alpha}{\partial \tau^2},\quad
  \alpha_{\delta\tau} = \frac{\partial^2 \alpha}{\partial \delta\, \partial \tau}
\]
(the crate's \texttt{HelmholtzDerivs} struct — see Part 3). Every residual
term type, including the non-analytic one, must supply all six.

\section{The derivative chain}

Differentiating Eq.~(3) from Part 1 term by term is mechanical but long; the
key structural fact that keeps it tractable is that every derivative can be
written as $\alpha_{r,i}$ times a \emph{ratio} function, which itself has a
clean recursive structure:
\begin{align}
  \frac{\partial \theta}{\partial \tau} &= -1, &
  \frac{\partial \theta}{\partial \delta} &= \frac{A_i}{\beta_i}
    \left[(\delta-1)^2\right]^{\frac{1}{2\beta_i}-1}(\delta - 1), \\
  \frac{\partial \Delta^{b_i}}{\partial \delta} &=
    b_i\, \Delta^{b_i - 1}\, \frac{\partial \Delta}{\partial \delta}, &
  \frac{\partial \Delta^{b_i}}{\partial \tau} &=
    -2\, b_i\, \theta\, \Delta^{b_i - 1}.
\end{align}
The full second-derivative expressions (implemented directly in
\texttt{eos.rs}, following CoolProp's \texttt{Helmholtz.cpp}) combine these
with the ordinary product rule applied to
$\alpha_{r,i} = n_i \delta\, \Delta^{b_i}\, \psi$. The pattern throughout is:
\emph{every} derivative is some polynomial combination of $\theta$, $\Delta$,
$\Delta^{b_i-1}$, $\Delta^{b_i-2}$, and the corresponding $\delta$/$\tau$
derivatives of $\theta$ and $\psi$ — nothing beyond elementary calculus, just
a long chain of it.

\section{The branch point: why $\delta=1,\ \tau=1$ is singular}

Look again at $\partial\theta/\partial\delta$ above. As $\delta \to 1$,
$(\delta-1)^2 \to 0$, and the exponent $\frac{1}{2\beta_i}-1$ is applied to
that vanishing quantity. For IAPWS-95 Water ($\beta_i = 0.3$), this
particular exponent is $\approx 0.67$ — positive, so this specific
expression is actually fine at $\delta=1$ (it evaluates to exactly $0$). The
\emph{second} $\delta$-derivative of $\theta$, needed for
$\alpha_{\delta\delta}$, instead uses the exponent $\frac{1}{2\beta_i}-2
\approx -0.33$ — \emph{negative} — so $\left[(\delta-1)^2\right]^{-0.33}$
diverges as $\delta \to 1$: $0^{-0.33} \to \infty$. A symmetric problem
appears in the $\Delta^{b_i-2}$ factor used by the second-$\delta$-derivative
of $\Delta^{b_i}$. In exact arithmetic, these divergences are known to
cancel against compensating factors elsewhere in the same expression (the
overall $\alpha_{r,i}$ and its first two derivatives remain finite at the
critical point) — but evaluating the raw intermediate quantities in floating
point \emph{at exactly} $\delta=1$ hits a literal
\texttt{0\textsuperscript{negative}} (effectively infinite) before those
cancellations can happen.

\section{The guard: offset by a few ULPs}

CoolProp's own resolution (\texttt{Helmholtz.cpp}, ported verbatim into this
crate's \texttt{accumulate\_non\_analytic}) is not to try to derive the
closed-form limit of every intermediate quantity — it is simpler and
numerically robust to just \emph{avoid} evaluating exactly at the branch
point:
\begin{lstlisting}[language=Rust]
let delta = if (delta - 1.0).abs() < 10.0 * f64::EPSILON {
    1.0 + 10.0 * f64::EPSILON
} else {
    delta
};
let tau = if (tau - 1.0).abs() < 10.0 * f64::EPSILON {
    1.0 + 10.0 * f64::EPSILON
} else {
    tau
};
\end{lstlisting}
Ten units in the last place ($10 \times$ \texttt{f64::EPSILON} $\approx
2.22\times10^{-15}$) is far smaller than any physically meaningful state
resolution, so this offset is invisible to any real calculation — but it is
large enough that IEEE~754 double-precision division and fractional
exponentiation near that point behave as ordinary (very large but finite)
numbers rather than literal infinities or NaNs, and the term's genuine
\emph{finite} limiting value falls out correctly from the surrounding
arithmetic.

\section{Why this is trustworthy: verification strategy}

A guard like this is easy to get subtly wrong (e.g.\ offsetting the wrong
variable, or by the wrong sign), so it is verified two ways, both covered in
Part 4: first, by evaluating \emph{exactly} at Water's tabulated critical
point $(T_c, \rho_c)$ — the offset guard fires on \emph{both} $\delta$ and
$\tau$ simultaneously, the most demanding case — and checking the resulting
pressure against IAPWS-95's defining value; second, by evaluating one degree
below $T_c$ on the critical isochore, where the guard does not fire at all,
as a monotonicity sanity check that the two code paths agree in the limit.

\printbibliography
\end{document}
